\chapter{Results}

\section{RNA-Seq workflow}
\subsection{Industrial RNA-Seq analysis pipeline}

\subsection{Automated gene noise filtering in RNASeq}

\section{Biological networks}

\subsection{Application to Covid19: unravel the pathways occurring in the \enquote{cytokine storm}}

\subsubsection{Motivation of the study}
\subsubsection{Results}
\subsubsection{Discussion}

\subsection{A local-score approach for the identification of key biological drivers}

\section{Model variability in biological systems}

\subsection{Gaussian Mixture Models in R}

\subsubsection{Motivation}
\subsubsection{Meta-analysis workflow and curation}
\subsubsection{Benchmark of \glsfmtshort{gmm}}
\paragraph{Methods}
\todo[fancyline]{integrate the numerical tip to deal with underflow + derivation of the overlap}
\paragraph{Univariate setting results}
\paragraph{Multivariate setting results}
\paragraph{Robustness}

\todo[fancyline]{outliers + skewed distribution}

\subsubsection{Discussion}


\subsection{Application to Sjögren's disease}
\subsubsection{Motivation of the study}
\subsubsection{Results}
\subsubsection{Discussion}


\subsection{Robust deconvolution of blood samples using the transcriptomic structure}

\subsubsection{Motivation of the study}

\subsubsection{Methods}

\paragraph{Online methods}

\paragraph{Overview of DECOVAR}
\subparagraph{Deconvolution model}

\subparagraph{Generative model}

\subparagraph{Derivation of the \glsentryshort{mle}}

\ref{pr:matrix-calculus} will be relevant for the derivation of the gradient:

\begin{Property}{Common matrix calculus formulas}{matrix-calculus}
Given three matrices $\boldsymbol{A}$, $\boldsymbol{C}$, $\boldsymbol{D}$ and two vectors of same size $\boldsymbol{b}$ and $\boldsymbol{e}$ not depending of variable $\boldsymbol{p}$, the partial derivative with respect to $\boldsymbol{p}$ of the quadratic form $(\boldsymbol{A} \boldsymbol{p} + \boldsymbol{b})^\top \boldsymbol{C} (\boldsymbol{D} \boldsymbol{p} + \boldsymbol{e})$ is given by \ref{eq:general-calculus-formula}:
\begin{equation}
\label{eq:general-calculus-formula}
    \frac{\partial (\boldsymbol{A} \boldsymbol{p} + \boldsymbol{b})^\top \boldsymbol{C} (\boldsymbol{D} \boldsymbol{p} + \boldsymbol{e})}{\partial \boldsymbol{p}} = (\boldsymbol{D} \boldsymbol{p} + \boldsymbol{e})^\top \boldsymbol{C}^\top \boldsymbol{A} + (\boldsymbol{A} \boldsymbol{p} + \boldsymbol{b})^\top \boldsymbol{C} \boldsymbol{D}
\end{equation}

Notably, using the transpose properties enumerated in \ref{transpose-properties}, setting $\boldsymbol{A} = \boldsymbol{D} = - \boldsymbol{x}$ as vectors living in $\mathrm{R}^G$, $p$ a real scalar and $\boldsymbol{b} = \boldsymbol{e} = \boldsymbol{y}$ the derivative \ref{eq:general-calculus-formula} and additionally coercing $\Theta$ to be symmetric yields:
\begin{equation*}
\label{eq:general-calculus-formula-easy}
    \frac{\partial (\boldsymbol{y} - \boldsymbol{x} p)^\top \Theta (\boldsymbol{y} - \boldsymbol{x} p)}{\partial p} = -2  (\boldsymbol{y} - \boldsymbol{x} p)^\top \Theta \boldsymbol{x} = -2 \boldsymbol{x}^\top \Theta (\boldsymbol{y} - \boldsymbol{x} p)
\end{equation*}

\footnote{Other matrix calculus formulas and notations are available on \href{https://en.wikipedia.org/wiki/Matrix_calculus\#Scalar-by-vector}{Matrix calculus}.}
\end{Property}

\begin{equation*}
\begin{array}{lcl}
\frac{\partial \ell_{\boldsymbol{y} | \boldsymbol{X}, \Sigma}}{\partial p_j}& =& \frac{\partial \log\left(\DET(\boldsymbol{\Theta})\right)}{\partial p_j} -\frac{1}{2} \left[\frac{\partial (\boldsymbol{y} - \boldsymbol{X} \boldsymbol{p})^\top}{\partial p_j}\boldsymbol{\Theta}(\boldsymbol{y} - \boldsymbol{X} \boldsymbol{p}) + (\boldsymbol{y} - \boldsymbol{X} \boldsymbol{p})^\top\frac{\partial\boldsymbol{\Theta}}{\partial p_j}(\boldsymbol{y} - \boldsymbol{X} \boldsymbol{p}) + (\boldsymbol{y} - \boldsymbol{X} \boldsymbol{p})^\top\boldsymbol{\Theta} \frac{\partial (\boldsymbol{y} - \boldsymbol{X} \boldsymbol{p})}{\partial p_j} \right]\\

&=&-\Tr \left(\boldsymbol{\Sigma}^{-1} \frac{\partial \boldsymbol{\Sigma}}{\partial p_j} \right) - \frac{1}{2} \left[ - \boldsymbol{x_{.j}}^\top\boldsymbol{\Theta}(\boldsymbol{y} - \boldsymbol{X} \boldsymbol{p}) - (\boldsymbol{y} - \boldsymbol{X} \boldsymbol{p})^\top\Theta\frac{\partial \Sigma}{\partial p_j}\Theta(\boldsymbol{y} - \boldsymbol{X} \boldsymbol{p}) - (\boldsymbol{y} - \boldsymbol{X} \boldsymbol{p})^\top\boldsymbol{\Theta}  \boldsymbol{x_{.j}} \right] \\

&=& -2p_j \Tr \left(\boldsymbol{\Sigma}^{-1}\right) + (\boldsymbol{y} - \boldsymbol{X} \boldsymbol{p})^\top\boldsymbol{\Theta}  \boldsymbol{x_{.j}} \, + p_j (\boldsymbol{y} - \boldsymbol{X} \boldsymbol{p})^\top\Theta \Sigma_j \Theta (\boldsymbol{y} - \boldsymbol{X} \boldsymbol{p})
\end{array}
\end{equation*}

\subsubsection{Results}

\subsubsection{Discussion: main challenges of cellular deconvolution}

